% =========================================================================================
% Chapter 5 -- Perception Agents
%
% Each module is written as a complete investigation ending in a decision: objective, data,
% preprocessing, features/model, training, results, error analysis, calibration, and the
% configuration retained. Rejected experiments belong inside the module that prompted them.
% Cross-module synthesis is Chapter 8, not here.
% =========================================================================================
\section{Modeling: Perception and Triage Agents}
\label{sec:perception}

\subsection{General Architecture}
\label{subsec:perception-arch}

Four modules are described in this chapter: one over structured data and three over images.
They differ in task, model family and evaluation metric, but share four design decisions that
are stated once here so that each module section need only report what it does differently.

\paragraph{Grouped cross-validation by case.} Every split is made at the level of the patient
or the case, never the image or the clip. Three of the four datasets provide frame-level labels
only, so case identity had to be reconstructed before any model was trained. Splitting on
frames instead would place different views of the same patient on both sides of the split, and
every reported figure would measure memorisation as much as generalisation.

\paragraph{Calibrated confidence fitted out of fold.} Each module emits a probability, and each
probability is calibrated by a method chosen to suit the shape of its miscalibration. The
calibrator is always fitted on data the model did not train on and never on the test set. A
module whose calibrator is missing reports \texttt{confidence\_calibrated: false} rather than
emitting a raw score that looks like a probability.

\paragraph{An explicit declaration of scope.} Each module states what it cannot do: which
findings it does not model, whether it has a healthy class, and the domain on which it was
trained. These declarations are part of the output, not documentation, and they reach the
reasoning layer of Chapter~\ref{sec:reasoning} where they bound what may be concluded.

\paragraph{A common output format.} All four write through the schema of
Section~\ref{subsec:output-contract}, which is what allows a segmentation network, two
classifiers and a tree ensemble to be consumed by a single reasoning layer.

The reporting convention of Section~\ref{subsec:methodology} does most of its work in this
chapter: each module is presented as it was built, measured and concluded, and every
configuration that was tested and rejected is reported inside the module that prompted it
rather than omitted.

% =========================================================================================
\subsection{Triage Agent}
\label{subsec:triage-agent}

\subsubsection{Objective}

The Triage Agent estimates an urgency tier from structured patient data alone: vital signs,
age, sex, arrival mode and presenting complaint. Its output is a tier with a calibrated
confidence, and the calibration is what makes that confidence usable alongside an imaging
finding's.

What it is not for matters as much. It does not diagnose, it proposes no differential, and it
never sees an image. It answers one question --- how urgently does this patient need to be
seen --- and the decision that follows from that answer is taken elsewhere.

\subsubsection{Dataset}

The agent is trained on the pooled emergency dataset described in Chapter~\ref{sec:data},
combining three independent sources: NHAMCS, an Iranian emergency department series, and
MC-MED. Provenance and licensing are reported there and not repeated here.

The target is a three-tier urgency label --- low, medium, high --- rather than the five-level
scales used in hospital triage. The reason is architectural. The reasoning layer consumes the
tier to decide whether to escalate, not to assign a waiting time, and a three-way decision is
what that consumer needs. The five-level formulation was also trained, and
Section~\ref{subsec:triage-results} reports why it was not retained.

\subsubsection{Preprocessing}

The most consequential preprocessing step concerns sentinel values. Vital signs recorded as
zero in these datasets mean ``not measured'', not a physiological zero: a systolic blood
pressure of 0 in a discharged patient is a missing record, not a finding. Left as numbers they
would be read as extreme readings, and a model would learn that the most abnormal-looking
patients are the ones nobody measured.

Every such sentinel was converted to a genuine missing value. Gradient-boosted trees handle
missing values natively, learning a default direction at each split rather than requiring
imputation, which is part of why that model family was chosen: the alternative --- imputing a
plausible value --- would fabricate a measurement that was never taken, which is the same error
this project refuses elsewhere.

\subsubsection{Feature Engineering}

Two groups of features were added, encoding how a clinician reads a set of vitals rather than
leaving the model to rediscover those relationships from limited data. The shock index (heart
rate divided by systolic pressure) captures a relationship between two vitals that neither
carries alone. A partial qSOFA score was computed from the components available in these
datasets. Red-flag threshold indicators mark values that cross the bounds a triage nurse is
trained to react to.

More valuable than what was added is what was removed. Disposition-related fields --- whether
the patient was admitted, transferred or discharged --- were present in the source data and
produced a markedly higher accuracy when included. That accuracy was not clinical signal. The
disposition is decided after triage and partly because of it, so including it lets the outcome
leak into the input and the model learn to read the answer rather than predict it. Every figure
reported in this chapter is from the model without those fields.

A coarse grouping of the reason-for-visit codes was also tested, on the hypothesis that
collapsing hundreds of codes into a few clinical categories would generalise better across
sources. It lost accuracy and was rejected: the fine-grained codes carry real signal that the
grouping discards.

\subsubsection{Model Selection}

Gradient-boosted trees were compared against logistic regression, a random forest, a small
multilayer perceptron and CatBoost on identical features and identical splits. None beat the
gradient-boosted model.

That result licenses an interpretation that would otherwise be unavailable. When five model
families of different inductive biases converge on a similar ceiling, the ceiling is a property
of the data rather than a limitation of the chosen architecture. Section~\ref{subsec:triage-error}
reads the errors on that basis: they are treated as reflecting genuine ambiguity in the
labelling of urgency rather than as evidence that a better model was available and not found.

\subsubsection{Training}

Training used a temporal split: the model fits the earlier year and is tested on the later one.
This is stricter than a random split and matches how the model would actually be used, on
patients arriving after those it learned from.

Class weighting compensates for the imbalance between tiers, with additional weight on the most
urgent classes so that under-triage is penalised more heavily than over-triage. A more
aggressive weighting was tested on the reasoning that safety should dominate: it cost
substantial overall accuracy for negligible improvement in dangerous errors, and was rejected.
The retained weighting was found by measurement across a sweep rather than assumed.

\subsubsection{Results}
\label{subsec:triage-results}

On 16\,180 held-out encounters the three-tier model reaches an accuracy of 0.680 and a macro
F1 of 0.666, against a chance level of 0.333 for three classes. Macro F1 is the honest headline
because it weights the three tiers equally; plain accuracy is inflated by the largest tier.

Per-tier performance is uneven in a way a pooled figure conceals. The high-urgency tier is
recognised best (F1 0.784 on 6\,465 encounters), the medium tier next (0.620 on 6\,263), and
the low tier worst (0.594 on 3\,452). The model is most reliable exactly where the consequences
of error are greatest, which is the desirable direction, though it should not be over-read: the
high tier is also the best-represented.

Performance also varies substantially by source:

\begin{table}[H]
  \centering
  \small
  \caption{Triage accuracy by source. A single pooled figure conceals this spread.}
  \label{tab:triage-per-source}
  \begin{tabular}{@{}lc@{}}
    \toprule
    Source & Accuracy \\
    \midrule
    Iranian ED series & 0.789 \\
    MC-MED & 0.652 \\
    NHAMCS & 0.617 \\
    \bottomrule
  \end{tabular}
\end{table}

A spread of 17 percentage points across sources is the most important number in this section.
It says that the pooled figure is an average over populations the model handles differently,
and it is the closest thing to external validation any module in this project has.

The five-level formulation, trained on the same data, reached an exact-match accuracy of 0.547
with a macro F1 of 0.394 --- substantially worse, and worse in a specific way: the two smallest
classes (immediate, $n=139$; non-urgent, $n=307$) were recognised at F1 0.170 each, barely
above chance. The three-tier target was retained because it is what the reasoning layer
consumes and because the finer target could not be learned reliably from this data.

Context makes these numbers interpretable. Published agreement between trained triage nurses
assigning the same scale to the same patients is itself imperfect, and a model that approaches
the level at which human raters agree with each other sits in a different position from one
that simply scores low. The comparison bounds what can be claimed rather than excusing the
result: the ceiling here is substantially set by the consistency of the labels.

\subsubsection{Error Analysis}
\label{subsec:triage-error}

Not all errors cost the same. Predicting high for a low-urgency patient wastes resources;
predicting low for a high-urgency patient can cost a life. Overall accuracy weights these
identically, so it is the wrong metric to optimise on its own.

For an ordinal target, accuracy-within-one-level is the appropriate companion metric. On the
five-level formulation the model was within one acuity level on 93.4\% of encounters and off by
two or more on 6.6\%, and among the 1\,734 genuinely high-acuity cases it was within one level
on 85.6\%, leaving 14.4\% dangerously under-triaged. That figure --- not the headline accuracy
--- is what a clinical reader should be shown.

A safety threshold addresses it directly: a prediction of ``low'' is forbidden when the model's
probability for ``high'' exceeds a fixed bound, in which case the prediction is raised. Sweeping
that bound gives the trade-off in Table~\ref{tab:triage-threshold}.

\begin{table}[H]
  \centering
  \small
  \caption{Safety threshold sweep. A threshold of 1.00 means no override.}
  \label{tab:triage-threshold}
  \begin{tabular}{@{}ccc@{}}
    \toprule
    Threshold & Dangerous miss rate & Overall accuracy \\
    \midrule
    1.00 (none) & 3.5\% & 0.6824 \\
    0.30 & 2.6\% & 0.6814 \\
    \textbf{0.20 (retained)} & \textbf{1.5\%} & \textbf{0.6802} \\
    0.15 & 0.1\% & 0.6436 \\
    0.10 & 0.0\% & 0.6308 \\
    \bottomrule
  \end{tabular}
\end{table}

The retained threshold of 0.20 removes more than half the dangerous misses --- 3.5\% to 1.5\%
--- at a cost of 0.2 percentage points of overall accuracy. Below 0.20 the trade turns sharply:
0.15 nearly eliminates dangerous misses but costs almost four points of accuracy, and by 0.10
the model is overriding so often that it has stopped discriminating. The knee of that curve was
found by measurement, and 0.20 sits at it.

\subsubsection{A Sanity Check on Constructed Cases}

Alongside the test-set metrics, five cases were written by hand to span the acuity spectrum and
the model was asked to grade them. The purpose is not measurement --- five constructed cases
support no statistical claim --- but a check that the model moves in the direction a clinician
would expect as the presentation changes, which an aggregate accuracy cannot show.

\begin{table}[H]
  \centering
  \small
  \begin{tabularx}{\linewidth}{Xllr}
    \toprule
    \textbf{Constructed case} & \textbf{Expected} & \textbf{Predicted} & \textbf{Conf.} \\
    \midrule
    Healthy young adult, minor ankle sprain, walk-in & low & low & 52.2\% \\
    Middle-aged, moderate abdominal pain, stable vitals, walk-in & medium & medium & 47.8\% \\
    Elderly, tachycardic, hypotensive and hypoxic, ambulance & high & high & 83.2\% \\
    \addlinespace[2pt]
    Elderly, ambulance, severe chest pain, tachycardic & high & \emph{medium} & 61.9\% \\
    Young adult, high fever and tachycardic, walk-in & medium & \emph{high} & 37.9\% \\
    \bottomrule
  \end{tabularx}
  \caption{Five hand-constructed cases. Three recover the expected tier.}
  \label{tab:triage-sanity}
\end{table}

\begin{figure}[H]
  \centering
  \includegraphics[width=\linewidth]{fig_triage_sanity.png}
  \caption{Predicted tier probabilities for the cases of Tab.~\ref{tab:triage-sanity}. The
    septic-shock-like presentation, where several vital signs are simultaneously abnormal, is
    the only one graded with high confidence (0.832); the remaining four distribute their mass
    across two or three tiers. This is a behavioural check on constructed inputs, not a
    performance measure.}
  \label{fig:triage-sanity}
\end{figure}

Both failures are confusions between \emph{adjacent} tiers, and neither confuses low with high
in either direction. They are not equivalent, however. The febrile young adult is over-triaged
at 37.9\% confidence --- the harmless direction, and without conviction. The elderly patient
with severe chest pain is under-triaged at 61.9\%, which is the combination that matters: the
wrong direction, stated with some confidence.

That case is the reason the tier is never the sole determinant of urgency. Reaching the
reasoning layer as \emph{medium}, it can still be escalated by an imaging finding, a missing key
laboratory value, or disagreement between the agents.

\subsubsection{Calibration}

Confidence is calibrated by Platt scaling fitted on a held-out slice of the training data,
never on the test set. Expected calibration error falls from 0.044 to 0.036.

That is a real improvement and a small one, because the model was already reasonably
calibrated. It is worth stating plainly what calibration did and did not do: accuracy is
essentially unchanged. Calibration does not make a model more accurate; it makes its stated
confidence mean something, which is what allows the reasoning layer to weigh a triage tier
against an imaging finding on a comparable scale.

\subsubsection{Final Decision}

The retained configuration is a gradient-boosted tree ensemble over the three-tier target,
trained on a temporal split with moderate class weighting, Platt-calibrated, and guarded by a
safety threshold of 0.20 on the high-urgency probability.

The limitations declared to the reasoning layer are that the model sees no imaging, that its
accuracy varies by source from 0.617 to 0.789, and that it estimates urgency rather than
diagnosis.

\subsubsection{Two Models, and Why There Are Two}
\label{subsec:triage-deployment}

The model described above is fitted on the full manifest, and part of its signal comes from
fields a hospital registration system assigns: the coded reason-for-visit categories, and the
source and year of the record. A clinician entering a patient at the bedside supplies none of
those. The intake screen asks for a presenting complaint in free text.

Restricting the fitted model to what the intake does supply was measured rather than assumed,
and the result was decisive: accuracy fell from 0.680 to 0.495, and on the held-out split the
model stopped predicting \textsc{low} altogether. A model starved of the features it was fitted
with is not the model whose accuracy was reported, and deploying it that way would put a figure
in this report that no deployed prediction could reach.

\begin{keystatement}
The response was not to restrict the model but to fit a second one, on exactly the fields the
interface collects. Training features, interface fields and deployment inputs are the same set.
\end{keystatement}

That second model is trained by the same procedure as the first --- the same temporal split, the
same class weighting, the same calibration slice carved out of the training data, the same
safety threshold --- so the two differ in their features and in nothing else.

\begin{table}[H]
  \centering
  \small
  \begin{tabular}{lrr}
    \toprule
    & \textbf{Reference} & \textbf{Deployed} \\
    \midrule
    Macro F1                    & 0.666 & 0.567 \\
    Accuracy                    & 0.680 & 0.630 \\
    Expected calibration error  & 0.036 & \textbf{0.030} \\
    Recall, \textsc{high}       & 0.784 & 0.773 \\
    Recall, \textsc{medium}     & 0.595 & \textbf{0.700} \\
    Recall, \textsc{low}        & \textbf{0.638} & 0.232 \\
    True \textsc{high} graded \textsc{low} & 1.55\% & \textbf{1.36\%} \\
    \bottomrule
  \end{tabular}
  \caption{The reference model and the deployed model on the same held-out split of 16,180
    visits. The deployed model is fitted on 21 features derived from the ten the intake
    collects; the reference model adds the registration fields.}
  \label{tab:triage-two-models}
\end{table}

The headline metric is lower, and the loss is not spread evenly. Recall on \textsc{high} is
essentially preserved and the dangerous miss --- a genuinely high-acuity patient graded
\textsc{low} --- is slightly \emph{better} than the reference model's. Almost the whole cost
falls on \textsc{low}: without the registration fields the model cannot confidently identify
low-acuity patients and errs towards \textsc{medium} instead. For decision support that is the
safe direction to lose accuracy in, and it is worth saying so explicitly rather than letting a
reader infer it from a table.

\paragraph{The two are registered separately.} Offering the deployed model as a replacement for
the reference model was tried and refused by the promotion gate of
Section~\ref{subsec:mlops}, for a macro-F1 regression of 0.099 that was both true and beside
the point. The gate compares successive versions of one model on one task; these are two models
answering different questions --- what the classifier achieves given complete records, and what
it achieves given what is actually typed. They now occupy separate production lines with
identical rules and unloosened tolerances, so a future deployed model is compared against the
deployed model it would replace and never against an artefact it was never built to beat.

\subsubsection{Suggestion, Confirmation, and Who Decides}

The deployed model does not set the urgency tier. It proposes one, and the clinician confirms
or changes it; the tier that reaches the reasoning layer is whichever the clinician left in the
control.

The interface shows the proposal and the observations it rests on, in clinical terms --- the
suggested tier, and either \emph{based on all recorded observations} or a count and the names
of those not taken. It does not show the probability vector, the model's name, or its
evaluation figures. A screen reporting \emph{0.646 medium, macro F1 0.567} asks a reader to
audit a classifier when their task is to see a patient, and those figures belong in this report
rather than at the bedside.

What is recorded with the encounter is fuller than what is displayed: the suggested tier, the
probability distribution behind it, the tier the clinician settled on, which observations were
available, and whether the two agreed. The record is written by the server from the state at
the moment of assessment, not accepted from the request, because a record a caller can write is
not a record --- it could report an agreement that never happened.

\begin{keystatement}
When the two differ, the system records that they differ. It does not record that the model was
wrong, because it has no ground truth with which to decide that.
\end{keystatement}

\paragraph{A defect this exposed.} The urgency control did not exist until late in the project.
The intake form carried a fixed value of \textsc{medium} and sent it with every encounter, so
the escalation trigger that fires when triage and imaging disagree could not fire for any
patient entered through the interface. The rule was correct throughout and its own tests passed,
because those tests call the conflict detector directly with a \textsc{low} tier. A rule can be
correct and never run, and no test that exercises the rule can reveal that: it took following
the value from the screen to the reasoning layer. The end-to-end checks added with the control
traverse that path over HTTP, which is where the gap was.

\paragraph{A second defect, in the suggestion itself.} The first implementation read the
model's probability vector in clinical order --- low, medium, high. The label encoder sorts its
classes, so the columns are ordered high, low, medium, and a hypotensive, hypoxic, tachycardic
81-year-old arriving by ambulance was graded \textsc{low} urgency at 0.87. Nothing raised: the
shapes matched, the probabilities summed to one, and every field was populated. The mapping is
now read from the artefact rather than written as a literal, and a test compares two
constructed patients by \emph{order} rather than by label, so that it survives retraining and
still fails if the vector is ever inverted again.

% =========================================================================================
\subsection{Ultrasound Agent: Organ Routing}
\label{subsec:ultrasound-agent}

The Ultrasound Agent presents one entry point. A study arrives with an organ and the
appropriate inputs --- two frames for cardiac, one still for gallbladder, a list of frames for
lung --- and is dispatched to the module that handles it.

An organ with no implemented module is not omitted and is not answered with a low-confidence
guess. It returns \texttt{status: "not\_supported"} in the ordinary output format, with an
empty finding list and a scope note. Vascular and FAST are in this position.

The consequence is that the reasoning layer never branches on whether a key exists. An
unsupported organ and an assessed one differ in the value of a field, not in the shape of the
object, and the distinction between ``examined and negative'' and ``never examined'' is carried
explicitly rather than inferred from an absence. Section~\ref{subsec:output-contract} describes
the contract this relies on.

\subsubsection{From notebook to importable module}
\label{subsubsec:agent-extraction}

The per-organ prediction code originally lived only in a Jupyter notebook. Nothing outside a
Colab session could execute it, and consequently no automated test could import it: the suite
proved that the output contract was \emph{enforced}, and could not prove that any perception
module \emph{honoured} it. That gap was reported as a limitation of this work.

It has since been closed for all three organs. The prediction code was extracted into an
importable module and is now exercised by sixteen tests grouped under a new safety property,
\emph{perception contract}, each running real inference on the shipped checkpoints. The lung
module was extracted first, and the cardiac and gallbladder modules followed once their
checkpoints and calibrators were retrieved from the training runs.

Two findings came out of the extraction, and both are more instructive than the extraction
itself.

\paragraph{The architecture had to be recovered from the checkpoints, not from the notebook.}
The two disagree. The notebook builds a two-layer classification head; every shipped checkpoint
contains a single linear layer of shape $4 \times 1280$. Loading the weights non-strictly
succeeds and leaves the mismatched layers at their random initialisation --- silently, and with
no error at any point. The architecture was therefore reconstructed from the checkpoint itself
and is now loaded strictly, so that any future divergence fails loudly instead of producing a
model that runs and means nothing.

\paragraph{The first test found a defect that had silenced the model.} The lung calibration
artefact described in Section~\ref{subsubsec:lung-calibration} was never exported from the
training notebook to a file. Both the notebook and the extracted module therefore fell back to
a decision threshold of $0.5$ for every finding. This model's outputs are compressed into
roughly the interval $[0.2, 0.8]$, so at $0.5$ almost nothing crosses: every study returned an
empty finding list, which the layers downstream read as a scan that was performed and showed
none of the modelled pathologies.

The failure is worth dwelling on because of its direction. A screening module that reports
false positives is annoying; one that silently reports every patient as clear is dangerous, and
it produces no error, no warning and no visible symptom. It was invisible precisely until the
code became importable and something asserted a property of it.

The tuned operating points had been written to disk by the training run all along, in the
per-split results files: $0.30$, $0.20$, $0.35$ and $0.45$ for B-lines, consolidation, pleural
effusion and pleural thickening respectively. The module now reads them back.

The correction was verified on LUS-BALD, an independent B-line localisation dataset that this
classifier never saw in training. Over 120 B-line-positive studies, the B-line output averages
$0.813$ and crosses its tuned threshold on every one of them, against an area under the ROC
curve of $0.919$ measured between those studies and random noise. Consolidation, correctly, sits
at $0.155$ on the same images --- and fires on none of them at a threshold of $0.5$. The model
was never the problem; the operating point was.

A distinction that the episode makes concrete, and that is easy to elide: a threshold states
\emph{where the decision boundary sits}, and calibration states \emph{what the number means}.
The thresholds were recovered. The Platt calibrators were not, and are genuinely lost, so the
deployed module continues to report \texttt{confidence\_calibrated: false} and emits a raw
sigmoid output. The boundary is tuned; the number is not.

\paragraph{Extracting the other two found three more defects, all of the same kind.} Cardiac and
gallbladder inference was extracted once their checkpoints were retrieved. The gallbladder
module transferred without incident. The cardiac module did not, and each of its three faults
produced a plausible output rather than an error --- which is why none had been noticed.

The first is a matter of arity. Ejection fraction is a \emph{comparison} between end-diastole
and end-systole, so a single still cannot produce one. Handed one frame as both, the arithmetic
computes a fractional change of zero and reports a normal ventricle. The interface now requires
both frames and refuses with a stated reason rather than returning a number.

The second is the serious one, and it was observed on real echocardiographic frames rather than
hypothesised. The systolic area was not checked for zero. A frame in which the segmentation
found no ventricle therefore entered the formula as an area of $0$, which the expression reads
as a ventricle that has ejected \emph{all} of its volume: an ejection fraction of 100\,\%,
banded \emph{normal}, and reported at the calibrator's maximum confidence. The direction of that
error is what makes it serious --- a failed measurement became a reassuring one. A frame in
which the ventricle cannot be found is a failed measurement, not a measurement of zero, and the
module now requires both areas to be non-zero, requires at least half the augmented views to
yield a usable pair, and treats a derived ejection fraction above 85\,\% as evidence of a
collapsed systolic segmentation rather than of a hyperdynamic heart.

The third is that calling the router without an image raised an exception. Every other
degradation in this system arrives as a report carrying \texttt{status: "failed"}, and this one
now does too, so the reasoning layer never has to catch an exception in order to stay correct.

\paragraph{An in-domain check, now that the module can be run outside a notebook.} The cardiac
module was evaluated against CAMUS's own recorded ejection fraction over thirty patients, using
the dataset's designated end-diastole and end-systole volumes. Twenty-eight produced a reading,
with a mean absolute error of 6.3 percentage points and a median of 5.2 --- consistent with the
6.9 reported from the training run in Section~\ref{subsec:cardiac-module}. The two refusals were
the patients with the highest true ejection fractions in the sample, where the module's upward
bias carries the estimate past the implausibility bound. That bound was left where it was rather
than adjusted after seeing which patients it caught, since tuning a threshold against the cases
it fails on is how a result becomes optimistic without anyone intending it.

The same run makes the generalisation gap of Section~\ref{subsec:cardiac-module} concrete. On an
echocardiographic source outside CAMUS the module segments a left ventricle averaging 157 pixels
against a myocardium of 954 --- an implausible ratio --- and the guards above refuse to produce a
band from it. The refusal is the correct behaviour and is also a measurement: the module does not
transfer, and it now declines rather than reporting a number from a segmentation it should not
trust.

% =========================================================================================
\subsection{Lung Module}
\label{subsec:lung-module}

\subsubsection{Objective}

The lung module detects four findings in a lung ultrasound clip: B-lines, consolidation,
pleural effusion and pleural thickening.

It is deliberately not a disease classifier. It reports what is visible in the images, and the
inference from B-lines to pulmonary oedema, or from consolidation to pneumonia, belongs to the
reasoning layer, which has access to the vitals, laboratory results and history that such an
inference requires. A module that reported ``pneumonia'' would be making a clinical judgement
from imaging alone.

\subsubsection{Dataset and the Multi-Label Formulation}

The four findings co-occur: a clip can show B-lines and a consolidation at the same time, and
frequently does. The task is therefore multi-label --- four independent probabilities --- rather
than single-label with one distribution over four mutually exclusive classes.

This is the first structural difference from the cardiac and gallbladder modules and it
propagates. It is the reason the shared schema carries a \emph{list} of findings rather than
one, and the reason the calibration section below fits four separate calibrators.

\subsubsection{Preprocessing and Clip Handling}

Eight evenly spaced frames are taken from each clip rather than every frame. Neighbouring frames
in an ultrasound clip are near-duplicates, so decoding all of them multiplies computation
without adding information. Frames are decoded once and cached.

Inference runs per frame and the frame probabilities are aggregated to a clip probability by
the \emph{median}. The choice matters and is not arbitrary. Frames from one clip are not
independent evidence about the patient; they are repeated looks at the same anatomy. The mean
would let one blurred or off-plane frame with an extreme probability drag the clip result, while
the median reports what the majority of the clip shows.

\subsubsection{Partially Labelled External Data}

Additional images were available that are labelled positive for B-lines only. They carry no
negatives and no labels at all for the other three findings. Treating the missing labels as
negatives would teach the model that these images show no consolidation, no effusion and no
thickening, none of which is known.

The solution is a per-sample mask applied to the loss, restricting each sample's gradient to
the columns whose labels are actually known. A B-lines-only image contributes to the B-lines
output and to nothing else.

The discipline around this matters as much as the mechanism: these samples entered the
\emph{training} set only. They never appear in validation, so no reported figure is computed
over data with partially invented labels.

\subsubsection{Model Selection and the Frozen-Backbone Discovery}

The single most valuable diagnostic in this project came from a configuration that was not
working. An earlier setup froze the entire EfficientNet-B0 backbone and trained only the
classifier head, on the standard reasoning that a small dataset cannot support fine-tuning a
full network.

That reasoning fails for ultrasound. Frozen ImageNet features encode edges, textures and object
parts learned from natural photographs, and ultrasound speckle resembles none of them. A frozen
backbone cannot adapt its representation, so the trainable head was being asked to separate
findings using features that did not describe them. Unfreezing the last blocks with a low
learning rate produced the largest single gain recorded anywhere in this project.

The general lesson is worth more than the specific gain: transfer learning assumes the source
and target domains share low-level structure, and that assumption should be checked rather than
inherited.

\subsubsection{Training}

Training used mixed precision, a cosine learning-rate schedule, and early stopping on validation
AUROC with a patience of three epochs against a budget of fifteen. Most folds stopped early. The
schedule is unremarkable and is reported for reproducibility rather than because anything was
learned from it.

\subsubsection{Results}
\label{subsec:lung-results}

Evaluation is 3-fold \texttt{GroupKFold} by case over 187 clips from 165 independent cases, with
predictions pooled out of fold and thresholds tuned on the folds not being evaluated.

\begin{table}[H]
  \centering
  \small
  \caption{Lung findings, pooled out-of-fold over 187 clips from 165 cases. ``pos'' is the
    number of positive clips available for that finding.}
  \label{tab:lung-results}
  \begin{tabular}{@{}lccccc@{}}
    \toprule
    Finding & AUROC & AP & Sensitivity & ECE & pos \\
    \midrule
    B-lines & 0.830 & 0.695 & 0.841 & 0.169 & 63 \\
    Consolidation & 0.832 & 0.726 & 0.600 & 0.146 & 60 \\
    Pleural effusion & 0.718 & 0.212 & 0.750 & 0.158 & 24 \\
    Pleural thickening & 0.773 & 0.453 & 0.652 & 0.195 & 46 \\
    \midrule
    \textbf{Macro} & \textbf{0.788} & & & & \\
    \bottomrule
  \end{tabular}
\end{table}

An earlier version of this evaluation reported AUROC between 0.745 and 0.808 and looked
comparable. It was inflated twice over: it split by clip rather than by case, so clips from the
same patient appeared on both sides of the split, and it tuned decision thresholds on the same
data it evaluated them on. Both inflations were removed and the figures above are the
defensible ones. They are reported here rather than the more flattering earlier numbers, and
the earlier numbers are reported alongside so that the difference is visible.

Pleural effusion is marked unreliable. Its 24 positive clips are too few to support a stable
threshold, and its average precision of 0.212 --- against an AUROC of 0.718 --- shows the
instability directly: the ranking is better than chance while the operating point is poor. That
mark is carried through the output contract to the reasoning layer, which discounts the finding
rather than weighing it like the others.

Pneumothorax is not modelled at all. The dataset contains zero positive cases, so no threshold
could be set and no confidence would mean anything. This absence is declared as a scope
limitation, and it is why an all-negative lung scan escalates in
Section~\ref{subsec:reasoning-escalation} rather than reading as reassurance.

\subsubsection{Error Analysis}

The two findings the model handles best, B-lines and consolidation, are also the two with the
most positive cases (63 and 60). The two it handles worst, pleural thickening and pleural
effusion, have 46 and 24. The ordering of performance follows the ordering of positive-case
counts exactly, which is the most parsimonious explanation available and does not require
appeal to intrinsic difficulty.

The clinically relevant confusions are between findings that co-occur: effusion and
consolidation appear together in the same clinical presentations, and a clip showing both gives
the model correlated evidence for each. Because the formulation is multi-label the model is not
forced to choose between them, which is the correct behaviour, but it also means the four
outputs are not independent and per-finding figures should not be multiplied together as if
they were.

\subsubsection{Calibration}
\label{subsubsec:lung-calibration}

Each finding receives its own Platt-scaled calibrator, fitted out of fold. One calibrator for
all four would be wrong: the four heads are miscalibrated to different degrees, as the ECE
column of Table~\ref{tab:lung-results} shows, ranging from 0.146 to 0.195.

Decision thresholds are tuned on the folds not being evaluated and applied to the held-out one.
Tuning and evaluating a threshold on the same data reports the best case rather than the
expected case, and was one of the two inflations removed from the earlier evaluation.

One limitation of the calibration is stated rather than hidden: it is fitted and scored on the
same out-of-fold predictions, so the reported ECE is optimistic.

A second limitation is procedural rather than statistical, and cost more. The calibrators were
fitted inside the training notebook and never written to a file, so nothing downstream could
load them. The consequences of that omission, and the recovery of the tuned thresholds from the
per-split results files, are described in Section~\ref{subsubsec:agent-extraction}. The
thresholds were retrievable; the calibrators were not. An artefact that exists only in the
memory of a notebook session does not exist.

\paragraph{Whether they could simply be refitted.} The obvious remedy is to regenerate the
out-of-fold predictions and fit the calibrators again. That was examined and rejected, on the
data rather than on effort. Calibration must be fitted on cases held out from both training
and threshold tuning, or the resulting probabilities are fitted to the same evidence they are
later scored against. Carving such a split out of this dataset leaves too little to fit
anything on:

\begin{table}[H]
  \centering
  \small
  \begin{tabular}{lrrr}
    \toprule
    \textbf{Finding} & \textbf{Positive clips} & \textbf{Positive cases} &
    \textbf{In a 20\% split} \\
    \midrule
    B-lines             & 63 & 53 & $\approx$ 11 \\
    Consolidation       & 60 & 46 & $\approx$ 9 \\
    Pleural thickening  & 46 & 39 & $\approx$ 8 \\
    \textbf{Pleural effusion} & \textbf{24} & \textbf{16} & $\approx$ \textbf{3} \\
    \bottomrule
  \end{tabular}
  \caption{Case-level positives per lung finding, and what a held-out calibration split would
    contain. The unit is the case, not the clip, for the reason established in
    Chapter~\ref{sec:data}.}
  \label{tab:lung-calibration-budget}
\end{table}

A Platt scaling fitted on three positive cases is not a calibration; it is a curve through
noise that would then be printed beside a clinical finding as a probability. The decision was
therefore to leave the module uncalibrated and to say so everywhere the number appears: the
report declares \texttt{confidence\_calibrated: false}, the scope string states that the
confidence is a raw sigmoid output and that the decision boundary is tuned while the number is
not, and \texttt{/api/health} reports the module as \emph{ready, uncalibrated}.

\begin{keystatement}
Reporting a raw score labelled as raw is honest. Manufacturing a calibration from three
positive cases and labelling it calibrated is not, and the label is what a clinician would
act on.
\end{keystatement}

This is the one place where the lung module is deliberately weaker than the cardiac and
gallbladder modules, which do carry fitted calibrators. It is a consequence of the dataset
size, not of the method, and it is the first thing additional lung data would fix.

\subsubsection{Final Decision}

The retained configuration is EfficientNet-B0 with the last blocks unfrozen, multi-label over
four findings, median aggregation across eight frames per clip, per-finding Platt calibration
and per-finding thresholds tuned out of fold.

The declared limitations are that pneumothorax is not modelled and cannot be excluded, that
pleural effusion is unreliable, that the module has no healthy class so an empty finding list
never means a normal study, and that the data are curated public clips rather than bedside
acquisitions.

% =========================================================================================
\subsection{Cardiac Module}
\label{subsec:cardiac-module}

\subsubsection{Objective}

The cardiac module segments the left ventricle, myocardium and left atrium in an
echocardiographic view, and derives an ejection fraction from the segmentation at end-diastole
and end-systole. The finding it reports is an EF band --- normal, reduced or severe --- with a
calibrated confidence.

It is the only pixel-level model in the system. Every other module classifies; this one labels
each pixel and computes a clinical quantity from the result, which makes its error analysis
different in kind from the others.

\subsubsection{Dataset}

Training uses CAMUS, which provides explicit patient structure and expert annotations, split by
patient. An additional external collection is used for validation only and never for training.
Provenance is in Chapter~\ref{sec:data}.

In-domain results are read against the published nnU-Net figures on CAMUS discussed in
Section~\ref{subsec:sota-existing}: Dice of 0.93 for the left ventricle, 0.86 for the myocardium
and 0.89 for the left atrium.

\subsubsection{Model and Loss}

The architecture is a U-Net with an EfficientNet-B0 encoder, 6.25 million parameters.

The loss combines cross-entropy with soft Dice, and the combination is necessary rather than
conventional. The structures of interest occupy a small fraction of the image, so a model
minimising cross-entropy alone can achieve a low loss by predicting background almost
everywhere: the pixel-wise objective is dominated by the majority class. Soft Dice measures
overlap between the predicted and true regions directly, and is insensitive to the size of the
background, so it rewards exactly what cross-entropy neglects.

Partially labelled masks are handled with an ignore index, so pixels whose true class is
unknown contribute no gradient rather than being treated as background.

\subsubsection{Training}

Early stopping monitors validation Dice, and the \emph{best} checkpoint is retained rather than
the last. The distinction matters when validation performance is non-monotone: keeping the last
checkpoint would report a model that had already begun to overfit.

\subsubsection{In-Domain Results}
\label{subsec:cardiac-results}

On held-out CAMUS patients the module reaches Dice of 0.928 for the left ventricle, 0.837 for
the myocardium and 0.875 for the left atrium, against the benchmark figures of 0.93, 0.86 and
0.89.

The left ventricle is essentially at benchmark. The myocardium and left atrium sit slightly
below. The left atrium is the weakest structure here, as it is in the published results ---
it has less consistent boundary contrast than the ventricle and its extent depends on the
imaging plane --- and reporting it as the weakest rather than averaging the three into one
number is what makes the comparison meaningful.

\subsubsection{External Validation and the Generalisation Gap}

On the external collection, left-ventricle Dice falls from 0.928 to 0.534. The drop is large
and is reported with both numbers rather than as a percentage.

The investigation proceeded in three steps, and the order is worth preserving because the
second step overturned the first.

\begin{enumerate}
  \item \textbf{A hypothesis formed from a handful of images.} Inspecting a small number of
        external cases suggested a labelling-convention mismatch: the external masks appeared
        to include or exclude structures at the boundary differently from CAMUS, which would
        depress Dice without the model having failed to find the ventricle.
  \item \textbf{Connected-component analysis over the full dataset, which disproved it.}
        Measuring mask fragmentation across every external case gave 2.4\% --- the predictions
        are coherent single regions, not scattered fragments, and the geometry does not show
        the systematic boundary offset the hypothesis predicted. The hypothesis was formed from
        a sample too small to support it and was abandoned.
  \item \textbf{What remains.} The gap is a genuine domain shift: different acquisition
        equipment, different operators, different imaging conventions. What it is \emph{not} is
        a labelling artefact, and it is not the model failing to localise the ventricle at all.
\end{enumerate}

The reversal is reported as method rather than concealed as an embarrassment. A conclusion
drawn from a handful of images and overturned by measurement over the full dataset is the same
error, in miniature, that the gallbladder module also produced
(Section~\ref{subsec:gallbladder-module}), and the recurrence is why it is stated twice.

Fig.~\ref{fig:cardiac-external-grid} shows what that failure looks like on individual frames,
and the visual form matters as much as the number. The predicted regions are single, smooth and
correctly placed over the ventricle; they are simply too small, cut off well inside the true
boundary. That is the signature of a decision threshold displaced by domain shift, not of a
model searching in the wrong part of the image.

\begin{figure}[H]
  \centering
  \includegraphics[width=\linewidth]{fig_cardiac_external_grid.png}
  \caption{Five frames drawn at random from the external collection: the image, the reference
    left-ventricle mask, and the prediction with its own Dice. The five values, 0.49 to 0.66,
    bracket the collection mean of 0.534 and are therefore representative rather than selected.
    Every prediction is one coherent region in approximately the right place and every one
    under-segments, which is the same conclusion the fragmentation measurement reached and the
    opposite of what a localisation failure would look like.}
  \label{fig:cardiac-external-grid}
\end{figure}

\subsubsection{Ejection Fraction}

Ejection fraction derived directly from segmented \emph{areas} carried a systematic bias: mean
absolute error of 12.2 percentage points against the reference values, with a correlation of
0.788.

The explanation is geometric rather than statistical. Ejection fraction is defined over
volumes, and area does not scale with volume: for a roughly ellipsoidal chamber, volume grows
as area to the power of $3/2$. Computing a fractional change from areas therefore understates
the fractional change in volume, systematically and predictably.

Applying the correction $\mathrm{EF} = \bigl(1 - r^{3/2}\bigr) \times 100$, where $r$ is the
ratio of end-systolic to end-diastolic area, reduced the mean absolute error from 12.2 to 6.9
percentage points --- a bias of 5.3 points removed --- and raised the correlation slightly to
0.795.

The result here is not the improved number. It is the agreement between the bias predicted from
the geometry and the bias actually observed. A correction derived from theory and then confirmed
by measurement is worth considerably more than a constant tuned until the error dropped, because
it will continue to hold on data this project has not seen.

On the three-band classification the module reaches a balanced accuracy of 63.4\% against a
chance level of 33.3\%. Plain accuracy is 65.0\% against a majority baseline of 66.0\%, and the
comparison shows why plain accuracy is the wrong headline: 66 of 100 patients fall in one band,
so predicting that band for everyone scores higher than the model while discriminating nothing.

Per-band recall is 0.727 for severe ($n=11$), 0.697 for reduced ($n=66$) and 0.478 for normal
($n=23$). The clinically important figure is the last line of the confusion matrix: \textbf{zero
of 100 cases confused a severe ventricle with a normal one}. The dangerous error --- reading a
failing heart as healthy --- did not occur in this cohort, though with 11 severe cases that is
a bound rather than a guarantee.

\subsubsection{A Small Error at a Cutoff is a Category Error}
\label{subsec:cardiac-boundary}

A mean absolute error of 6.9 points reads as a good result, and on its own it is. But the number
is not what a clinician acts on --- the band is, and the bands are separated by cutoffs at 30
and 55. An error far smaller than the reported average can move a patient across one.

The effect is visible on individual cases. Estimating a reference ejection fraction of 29 as 31
is a numeric error of 2.0 points, well inside the module's typical accuracy, and it changes the
clinical statement from \emph{severe} to \emph{moderate} dysfunction. The same arithmetic
applies at the upper cutoff, where an error of 3.7 points moved a reference value of 54 to 57.7
and turned a mildly reduced ventricle into a normal one.

The module therefore reports, with every ejection fraction, how far that value sits from the
nearest cutoff and whether it lies within the module's own measured error of one.

\begin{keystatement}
The margin is the module's held-out mean absolute error, not a chosen number. It is an
engineering sensitivity zone: it says this model cannot resolve which side of the cutoff the
patient is on, not that a cardiologist would call the value borderline.
\end{keystatement}

The distinction matters because a margin invented for the purpose --- five points, three points
--- would be presented to a clinician as though it carried clinical meaning. Deriving it from
the measured error makes it a statement about the instrument, which is what it is.

The consequence is visible and worth stating plainly: because 6.9 points is wide relative to a
\emph{reduced} band only 25 points across, a substantial share of predictions fall inside the
zone and are flagged. That is the honest reading of an error of this size, and a narrower margin
chosen to reduce the flag rate would have missed both of the band crossings described above.
This is also the reason balanced accuracy over the three bands, at 63.4\%, is a far less
flattering headline than the 6.9-point error --- and the more truthful one.

\subsubsection{Rejected Adaptations}

Six strategies were tried against the external gap and four were rejected.

\begin{table}[H]
  \centering
  \small
  \caption{Adaptations tested against the cardiac generalisation gap.}
  \label{tab:cardiac-rejected}
  \begin{tabularx}{\textwidth}{@{}lXl@{}}
    \toprule
    Strategy & Hypothesis and outcome & Decision \\
    \midrule
    Augmentation upgrade & Stronger augmentation would cover the domain shift; the measured
      difference was within run-to-run variation & rejected \\
    External integration (2 ways) & Adding external data to training would bridge the gap; both
      integrations degraded in-domain performance without closing it & rejected \\
    Fine-tuning (2 ways) & Fine-tuning on the external set would adapt the representation; it
      improved the external number substantially while degrading CAMUS performance to an
      unacceptable degree & rejected \\
    Volume correction & Area-derived EF is biased by a predictable geometric factor; MAE
      12.2\,pp $\rightarrow$ 6.9\,pp & \textbf{retained} \\
    \bottomrule
  \end{tabularx}
\end{table}

Fine-tuning is the instructive rejection. It worked, in the narrow sense that the external
number improved substantially. It was rejected because the improvement was purchased with
in-domain performance, and the in-domain distribution is the one the module is declared to
serve. A model that performs moderately everywhere and well nowhere is not obviously preferable
to one that performs well within a stated domain and declares that domain honestly.

The non-fine-tuned model was retained and its domain limitation is declared to the reasoning
layer rather than hidden by an averaged figure.

\subsubsection{Calibration}

Confidence is the fraction of eight fixed test-time augmentations that agree on the reported
band --- fixed rather than random, so the same clip always yields the same number --- and that
agreement fraction is calibrated by isotonic regression.

Isotonic regression suits this case because an agreement fraction is not a softmax output and
its distortion is not of a form a single temperature parameter can correct; a non-parametric
monotone fit can. Expected calibration error falls from 0.314 to 0.060, by far the largest
calibration correction of any module, which reflects how poorly a raw agreement fraction
behaves as a probability.

The isotonic fit also implies a confidence ceiling: the calibrator maps even unanimous
agreement to a value below 1.0, and that ceiling is declared in the output so that no
downstream component treats the module as certain.

\subsubsection{Final Decision}

The retained model is the CAMUS-trained EfficientNet-B0 U-Net without fine-tuning, with
volume-corrected ejection fraction reported as a three-band finding and isotonic-calibrated
confidence.

The declared scope is explicit: CAMUS-like four-chamber views, ejection fraction as an area
proxy rather than a volumetric measurement, and a documented drop on out-of-domain
acquisitions. What this forbids the reasoning layer from concluding is that the module's
confidence is meaningful on an image unlike its training domain --- and because the module
cannot tell whether a given image is in domain, that judgement is left to the clinician, with
the limitation stated in every report.

% =========================================================================================
\subsection{Gallbladder Module}
\label{subsec:gallbladder-module}

\subsubsection{Objective}

Single-label classification of gallbladder pathology at case level, over five reported classes:
cholelithiasis, acute cholecystitis, polyps and adenomyomatosis, carcinoma, and wall thickening.

\subsubsection{Dataset and the Mislabelled Class}

The dataset comprises 4\,826 frames from 199 cases, organised as nine folders by pathology.

A per-case inspection found that one folder's contents did not match its label. The
methodological history of that finding is more valuable than the finding itself. The class had
first been judged from a handful of frames drawn from a small number of cases, and looked
consistent. Sampling one frame per case across all cases in the folder reversed the conclusion
immediately.

The rule this establishes generalises well beyond this dataset: \textbf{frames from one case are
not a sample of a class}. A handful of frames from two cases shows what those two cases look
like, and any impression formed from them describes the cases rather than the category. The same
error appeared independently in the cardiac module, where a hypothesis formed from a few images
was overturned by measurement over the full dataset.

The affected folder was excluded, reducing the taxonomy from nine classes to eight for training.

\subsubsection{Taxonomy Selection}

Several groupings of the eight training classes into a reported taxonomy were plausible. Training
a model per candidate would have cost GPU time proportional to the number of candidates.

Instead, the confusion matrix already measured from the fine-grained model was collapsed
arithmetically according to each candidate grouping, estimating the accuracy each would achieve
without training anything. The candidate that this predicted would perform best was the one
taken forward, and it predicted 63.0\% balanced accuracy for collapsing the fine-grained model's
predictions.

The measured result after training the chosen configuration was 59.7\%. The agreement between a
prediction made from an existing confusion matrix and a figure measured after training is itself
a result: it justifies having spent GPU time on one configuration rather than several, and the
three-point shortfall marks the limit of the approximation rather than invalidating it.

\subsubsection{Training}

The model trains on the eight fine-grained labels and reports five, with class weights computed
per fold from the training portion only --- computing them over the whole dataset would leak the
test distribution into training.

The marginalisation from eight classes to five sums group probabilities rather than mapping the
argmax. The distinction is not cosmetic. A case whose probability mass is split across
cholecystitis, gangrenous cholecystitis and perforation --- three members of one reported group
--- has a high total probability of acute inflammation and no single dominant fine-grained class.
Mapping the argmax reports whichever member happens to be largest, discarding the agreement
between them; summing reports what the model actually believes about the group. Measured
directly, summing scored 63.0\% against 57.2\% for training on the merged labels.

\subsubsection{Results}
\label{subsec:gallbladder-results}

Evaluation is 5-fold \texttt{GroupKFold} by case, case-level predictions, over 199 cases with
frames capped at 25 per case so that longer studies do not dominate.

Balanced accuracy on the five reported classes is 59.7\% against a chance level of 20.0\%; plain
accuracy is 61.3\%. Per-class recall is given in Table~\ref{tab:gb-results}.

\begin{table}[H]
  \centering
  \small
  \caption{Gallbladder per-class recall, case level, with the one-standard-error bound at each
    class's case count.}
  \label{tab:gb-results}
  \begin{tabular}{@{}lcc@{}}
    \toprule
    Class & Cases & Recall \\
    \midrule
    Acute cholecystitis (any severity) & 66 & 0.70 $\pm$ 0.06 \\
    Carcinoma & 27 & 0.67 $\pm$ 0.10 \\
    Polyps / adenomyomatosis & 49 & 0.59 $\pm$ 0.07 \\
    Wall thickening & 25 & 0.56 $\pm$ 0.10 \\
    Cholelithiasis & 32 & 0.47 $\pm$ 0.09 \\
    \bottomrule
  \end{tabular}
\end{table}

The progression that produced this figure is worth reporting in full, because the improvement is
only interpretable against its starting point:

\begin{table}[H]
  \centering
  \small
  \caption{Balanced accuracy on the same 199 cases and the same reporting taxonomy.}
  \label{tab:gb-progression}
  \begin{tabular}{@{}lc@{}}
    \toprule
    Configuration & Balanced accuracy \\
    \midrule
    Original 9-class taxonomy, mislabelled folder included & 37.5\% \\
    Trained directly on the 5 merged classes & 57.2\% \\
    Collapsing the 9-class model's predictions & 63.0\% \\
    Retained: 8-class training, summed output, class-weighted & 59.7\% \\
    \bottomrule
  \end{tabular}
\end{table}

Grouping the five reported classes further into three clinical categories gives 63.6\% balanced
accuracy against a majority baseline of 50.8\%. That grouping is supported by the imaging only
weakly: 18\% of the fine-grained model's confusions stay inside their group, meaning most errors
cross group boundaries and the categories are not capturing a distinction the model can see.
The five-class output was retained on that basis.

\subsubsection{Rejected Alternatives}

Two alternatives were tested and rejected, both with the numbers that decided them.

\textbf{Training directly on merged labels} scored 57.2\% balanced accuracy against 63.0\% for
training fine-grained and merging probabilities afterwards. Collapsing the labels before
training discards the distinction between group members, and with it the model's ability to
accumulate evidence across them.

\textbf{USCL initialisation} was tested on the strength of the published gain discussed in
Section~\ref{subsec:sota-existing}: ultrasound-specific self-supervised pretraining should
transfer better to ultrasound than ImageNet weights. On a fold-0 bake-off it reached 29.8\%
balanced accuracy against 46.6\% for the ImageNet-initialised EfficientNet-B0 --- not a marginal
difference. The published result was not reproduced here, and the ImageNet initialisation was
retained.

\subsubsection{Calibration}

Temperature scaling, with a fitted temperature of $T = 0.87$. Expected calibration error moves
from 0.025 to 0.024.

The temperature below 1.0 indicates the model was mildly \emph{under}-confident rather than
over-confident, which is unusual and worth noting: the class-weighted loss, which penalises
errors on rare classes more heavily, appears to have flattened the output distribution slightly.
The correction is small because there was little to correct.

\subsubsection{Final Decision}

The retained configuration is EfficientNet-B0 trained on eight fine-grained classes with
per-fold class weights, reporting five classes by summing group probabilities, temperature
calibrated at $T = 0.87$.

The declared limitation that matters most is the absence of a healthy class. The dataset
contains no normal gallbladders, so every prediction is a choice among pathologies and none of
them means ``no disease''. The module therefore cannot exclude anything, and the reasoning layer
is told so explicitly: a gallbladder report is evidence about \emph{which} pathology, never
evidence that there is none. The remaining declared limitations are the 199-case sample, the
burned-in annotation markers present in 15--50\% of frames, and the teaching-atlas provenance.

\subsubsection{Reporting the Classes That Were Not Chosen}
\label{subsec:gallbladder-alternatives}

A module that names one diagnosis and stops leaves a reader unable to distinguish two very
different situations: a diagnosis that was weighed against the alternatives and preferred, and
a diagnosis reported because the alternatives were never in the label set at all. The lung
module does not have this problem --- it reports all four findings it screens, positive or
negative --- but a single-label module does, and the gallbladder module is single-label.

It therefore reports the four classes it ranked below the one it chose, with their
probabilities and the margin between the first and the second. On a cholecystitis case the
report reads: acute cholecystitis at 0.541, then cholelithiasis 0.234, carcinoma 0.175, polyps
0.039 and wall thickening 0.011.

\begin{keystatement}
These are not negative results. The five probabilities sum to one and describe one mutually
exclusive choice, so a class ranked lower was never screened for and found absent --- it was
outranked.
\end{keystatement}

The distinction is enforced in the schema rather than left to a convention. The ranked classes
occupy their own field, separate from the list of findings a module screened and did not see,
and a report carrying both is rejected as invalid: a module either screens independent findings
or chooses among exclusive ones, and one claiming to do both cannot be interpreted --- a reader
could not tell whether a label absent from the findings had been ruled out or merely outranked.

\paragraph{They are shown to the clinician and withheld from the language model.} The ranked
alternatives reach the interface but never enter the clinical state, and that is the point of
keeping them in a separate field. Everything in the clinical state becomes enumerated evidence
that the reasoning layer may cite; an alternative is not an observation about the patient but a
diagnosis the module considered and did not make. Citing ``cholelithiasis 0.23'' as evidence
would be fabrication with a number attached, and the separation makes it unrepresentable rather
than merely discouraged.

The margin is reported for a related reason. A confidence of 0.541 says how sure the module is;
it does not say how close the runner-up was. A margin near zero means the module effectively
could not separate two diagnoses, which is a different warning from low confidence and is not
recoverable from the top-line number alone.

% =========================================================================================
\subsection{Explainability Across the Imaging Modules}
\label{subsec:explainability}

Grad-CAM applies to all three imaging modules and is therefore treated once here rather than
three times.

The methodological point is that the maps were \emph{measured} rather than inspected by eye.
Looking at a dozen heatmaps and finding them plausible is not evidence; it is an invitation to
confirmation bias, since a viewer who knows where the pathology is will see the map pointing at
it. Three statistics were computed over the full sets instead:

\begin{itemize}
  \item \textbf{centroid spread} --- how much the centre of attention moves between images. A
        map that always points at the same place is responding to position rather than content;
  \item \textbf{area above half-maximum} --- what fraction of the image carries at least half
        the peak activation, which separates concentrated attention from diffuse;
  \item \textbf{normalised entropy} --- the dispersion of the attention distribution, referenced
        against the values for uniform attention (entropy 1.0) and for a single concentrated
        region.
\end{itemize}

An explicit test for centre bias accompanies them, because a heatmap that points at the middle
of every image regardless of content is not an explanation --- and ultrasound images place the
region of interest near the centre often enough that centre bias would look like success.

What these statistics support is that attention is concentrated and content-dependent rather
than uniform or positional. What they cannot support is that the attention falls on the
anatomically correct structure, which would require expert annotation of the relevant structure
in each image, and this project has none. The gallbladder module in particular showed a
pronounced central prior consistent with either a genuine anatomical focus or an acquisition
convention, and nothing measured here separates the two.

Fig.~\ref{fig:lung-gradcam} shows the lung maps, one clip drawn at random from the positives of
each of the four findings. It is included to make the measured statistics legible, and it
requires one caveat that the panel labels do not carry themselves.

\begin{figure}[H]
  \centering
  \includegraphics[width=\linewidth]{fig_lung_gradcam.png}
  \caption{Grad-CAM for the lung module: one clip per finding, top row the frame and bottom row
    the attention. The \textsc{pred} label is the \emph{highest-scoring} of the four findings,
    not the module's decision --- the module is multi-label and reports each finding
    independently against its own threshold, so a clip may carry two findings while only one
    can be named here. These panels therefore say nothing about whether the module classified
    the clip correctly.}
  \label{fig:lung-gradcam}
\end{figure}

The maps vary in how concentrated they are, and the first panel is the least concentrated of
the four. That is consistent with the statistics rather than in tension with them: the claim
established is that attention is concentrated \emph{on average} and moves with content, not
that every individual map is tight. Attention reaching the edge of the sector in several panels
is the observation that most warrants caution here, since the sector boundary is an artefact of
the acquisition geometry and carries no clinical information --- a model attending to it would
be reading the machine rather than the patient. Nothing measured here settles that, which is
why the boundary below is stated as firmly as it is.

Grad-CAM is therefore reported as a plausibility check that the models are not obviously
responding to the wrong thing, and explicitly not as evidence of clinical validity.
Section~\ref{subsec:results-explainability} states the same boundary at system level.

% =========================================================================================
\subsection{Perception Agent Output Contract}
\label{subsec:output-contract}

Every module writes through one schema. This section is the hinge between
Chapters~\ref{sec:perception} and~\ref{sec:reasoning}: it is what makes a segmentation network,
two classifiers and a tree ensemble usable by a single reasoning layer without that layer
knowing which produced what.

A report carries the organ, a status, a list of detected findings each with a calibrated
confidence, a list of findings screened for and \emph{not} detected, and a reliability block
declaring the module's scope.

Four validation rules are enforced, each for a reason established by a failure:

\begin{itemize}
  \item \textbf{Findings are always a list.} Lung is multi-label and a single-finding field
        could not represent a clip showing B-lines and consolidation together. Making the list
        universal means the reasoning layer never special-cases by organ.

  \item \textbf{A status of ``ok'' with nothing recorded is rejected.} A report claiming an
        assessment took place while recording neither a detected nor a screened-negative finding
        asserts nothing and would be indistinguishable downstream from an organ never examined.

  \item \textbf{A confidence above the declared ceiling is rejected.} The cardiac calibrator maps
        even unanimous test-time agreement to a value below 1.0; a report exceeding that ceiling
        indicates the calibration was not applied.

  \item \textbf{A label cannot appear in both the detected and not-detected lists.} The two
        lists make contradictory claims about the same finding, and silently preferring one
        would decide a contradiction the module should not have produced.
\end{itemize}

\paragraph{Four states on the screen the clinician reads.} The contract carries three states;
the interface shows four, because one of the three splits in two once a reader is involved.

\begin{table}[H]
  \centering
  \small
  \begin{tabularx}{\linewidth}{lX}
    \toprule
    \textbf{State} & \textbf{What it asserts} \\
    \midrule
    Detected & The module saw the finding and can support the claim \\
    \addlinespace[2pt]
    Detected --- weak evidence & The module reports it, on evidence too thin to act on alone.
      Pleural effusion rests on 16 positive cases and carries this permanently \\
    \addlinespace[2pt]
    Screened, not detected & The module looked for it and did not see it. An informative
      negative \\
    \addlinespace[2pt]
    NOT assessed --- not modelled & The module does not cover this finding at all. It appears
      in the list with no confidence value, because no examination of it took place \\
    \bottomrule
  \end{tabularx}
  \caption{The four states a lung finding can occupy in the interface.}
  \label{tab:four-states}
\end{table}

The fourth state exists for pneumothorax specifically. It was in the intended label set, has
zero positive examples in the usable data, and was therefore dropped --- and it is among the
findings lung ultrasound is most often used to look for. A list showing B-lines, consolidation
and pleural thickening all negative, with pneumothorax simply absent from it, invites the one
inference the whole system is built to prevent: the scan did not mention it, so it is not
there. Giving it a row of its own, with a dash where a confidence would be, makes the silence
visible.

The distinction the contract exists to preserve is between three states, not two: a finding
detected, a finding screened for and not seen, and a finding never assessed. The first two are
observations; the third is a gap in the record. A module that screened four findings and saw
none reports an empty \texttt{findings} list beside a populated \texttt{not\_detected} list ---
an informative negative result --- while an organ never examined is absent from the report
entirely.

Collapsing that distinction was a real defect rather than a hypothetical one. The lung module
once represented ``nothing above threshold'' with a placeholder entry in the findings list.
Everything in that list is marked as detected downstream, so the \emph{absence} of findings was
read as a positive finding: it suppressed the escalation trigger that fires when no finding is
detected, and counted as weak supporting evidence in the case-quality grade. An all-negative
scan escalated less readily than a positive one, which is backwards, and on exactly the scans
where an unmodelled pathology such as pneumothorax is most likely to be the explanation. The
correction is the \texttt{not\_detected} list described above, and it was verified end to end.
